\documentclass[11pt,a4paper]{article}

%========================
% PACKAGES
%========================
\usepackage[utf8]{inputenc}
\usepackage[T1]{fontenc}
\usepackage[french]{babel}
\usepackage{lmodern}
\usepackage{microtype}
\usepackage{amsmath,amssymb,amsthm,mathtools}
\usepackage{geometry}
\usepackage{fancyhdr}
\usepackage{xcolor}
\usepackage{enumitem}
\usepackage{tcolorbox}
\usepackage{array}
\usepackage{booktabs}
\usepackage{hyperref}
\usepackage{tikz}

\geometry{margin=1.7cm}
\setlength{\parindent}{0pt}
\setlength{\parskip}{0.45em}

%========================
% COULEURS
%========================
\definecolor{bleuprepa}{RGB}{20,55,110}
\definecolor{rougeprepa}{RGB}{160,35,35}
\definecolor{vertprepa}{RGB}{20,105,70}
\definecolor{orangeprepa}{RGB}{180,100,20}
\definecolor{grisclair}{RGB}{245,247,250}

%========================
% EN-TÊTE
%========================
\pagestyle{fancy}
\fancyhf{}
\lhead{\textsc{Analyse -- Fiche de référence}}
\rhead{\textsc{Fonctions usuelles et dérivation}}
\cfoot{\thepage}

%========================
% BOÎTES
%========================
\tcbset{
  enhanced,
  breakable,
  sharp corners,
  boxrule=0.7pt,
  colback=white,
  fonttitle=\bfseries
}

\newtcolorbox{definitionbox}[1]{
  colframe=bleuprepa,
  colback=blue!3!white,
  title={Définition -- #1}
}

\newtcolorbox{theorembox}[1]{
  colframe=vertprepa,
  colback=green!3!white,
  title={Théorème / Propriété -- #1}
}

\newtcolorbox{methodbox}[1]{
  colframe=orangeprepa,
  colback=orange!5!white,
  title={Méthode -- #1}
}

\newtcolorbox{warningbox}[1]{
  colframe=rougeprepa,
  colback=red!4!white,
  title={Attention -- #1}
}

\newtcolorbox{proofbox}[1]{
  colframe=black,
  colback=grisclair,
  title={Démonstration -- #1}
}

%========================
% COMMANDES
%========================
\newcommand{\R}{\mathbb{R}}
\newcommand{\N}{\mathbb{N}}
\newcommand{\dd}{\,\mathrm{d}}
\newcommand{\e}{\mathrm{e}}
\newcommand{\Id}{\mathrm{Id}}
\newcommand{\ch}{\operatorname{ch}}
\newcommand{\sh}{\operatorname{sh}}
\newcommand{\thh}{\operatorname{th}}
\newcommand{\Arcsin}{\arcsin}
\newcommand{\Arccos}{\arccos}
\newcommand{\Arctan}{\arctan}

%========================
% DOCUMENT
%========================
\begin{document}

\begin{center}
    {\Huge \bfseries Fiche complète d'analyse}\\[0.3em]
    {\Large Fonctions usuelles, dérivées et démonstrations importantes}\\[0.4em]
    {\large Niveau MPSI -- PCSI -- MP}\\[0.5em]
    \rule{0.85\linewidth}{0.5pt}
\end{center}

\tableofcontents
\newpage

%====================================================
\section{Notions fondamentales sur les fonctions}

\begin{definitionbox}{Fonction réelle}
Une fonction réelle d'une variable réelle est une application
\[
f : D \subset \R \longrightarrow \R.
\]
Pour tout $x \in D$, le réel $f(x)$ est appelé l'image de $x$ par $f$.
\end{definitionbox}

\begin{definitionbox}{Domaine de définition}
Le domaine de définition de $f$ est l'ensemble des réels $x$ pour lesquels l'expression $f(x)$ a un sens.

On le note souvent $D_f$.
\end{definitionbox}

\begin{methodbox}{Déterminer un domaine}
Pour trouver $D_f$, on vérifie :
\begin{enumerate}[label=\arabic*)]
    \item les dénominateurs ne doivent pas être nuls ;
    \item les racines carrées imposent une quantité positive ou nulle ;
    \item les logarithmes imposent un argument strictement positif ;
    \item les arcsinus et arccosinus imposent un argument dans $[-1,1]$ ;
    \item les compositions doivent être cohérentes.
\end{enumerate}
\end{methodbox}

\begin{definitionbox}{Parité}
Soit $D$ symétrique par rapport à $0$.
\begin{itemize}
    \item $f$ est paire si
    \[
    \forall x\in D,\quad f(-x)=f(x).
    \]
    \item $f$ est impaire si
    \[
    \forall x\in D,\quad f(-x)=-f(x).
    \]
\end{itemize}
\end{definitionbox}

\begin{warningbox}{Parité}
Avant d'étudier la parité, il faut toujours vérifier que le domaine est symétrique par rapport à $0$.
\end{warningbox}

\begin{definitionbox}{Périodicité}
Une fonction $f$ est périodique s'il existe $T>0$ tel que
\[
\forall x\in D_f,\quad x+T\in D_f
\quad\text{et}\quad
f(x+T)=f(x).
\]
Un tel réel $T$ est une période de $f$.
\end{definitionbox}

%====================================================
\section{Limites : rappels essentiels}

\begin{definitionbox}{Limite finie en un point}
Soit $f : D\to\R$, $a\in\R$ et $\ell\in\R$.

On dit que
\[
\lim_{x\to a} f(x)=\ell
\]
si
\[
\forall \varepsilon>0,\ \exists \eta>0,\ \forall x\in D,\quad
0<|x-a|\leq \eta \Longrightarrow |f(x)-\ell|\leq \varepsilon.
\]
\end{definitionbox}

\begin{definitionbox}{Continuité}
Soit $a\in D_f$. La fonction $f$ est continue en $a$ si
\[
\lim_{x\to a} f(x)=f(a).
\]
Elle est continue sur un intervalle $I$ si elle est continue en tout point de $I$.
\end{definitionbox}

\begin{theorembox}{Caractérisation séquentielle}
Soit $f:D\to\R$, $a\in\R$ et $\ell\in\R$.

On a
\[
\lim_{x\to a} f(x)=\ell
\]
si et seulement si, pour toute suite $(x_n)$ d'éléments de $D\setminus\{a\}$ telle que $x_n\to a$, on a
\[
f(x_n)\to \ell.
\]
\end{theorembox}

\begin{methodbox}{Montrer qu'une limite n'existe pas}
Pour montrer que $f$ n'a pas de limite en $a$, on peut trouver deux suites $(u_n)$ et $(v_n)$ telles que
\[
u_n\to a,\qquad v_n\to a,
\]
mais
\[
f(u_n)\to \ell_1,\qquad f(v_n)\to \ell_2,
\]
avec
\[
\ell_1\neq \ell_2.
\]
\end{methodbox}

\begin{theorembox}{Théorème des gendarmes}
Si, au voisinage de $a$,
\[
g(x)\leq f(x)\leq h(x)
\]
et si
\[
\lim_{x\to a}g(x)=\lim_{x\to a}h(x)=\ell,
\]
alors
\[
\lim_{x\to a}f(x)=\ell.
\]
\end{theorembox}

%====================================================
\section{Dérivabilité}

\begin{definitionbox}{Taux d'accroissement}
Soit $f:I\to\R$ et $a,x\in I$ avec $x\neq a$.

Le taux d'accroissement de $f$ entre $a$ et $x$ est
\[
\frac{f(x)-f(a)}{x-a}.
\]
\end{definitionbox}

\begin{definitionbox}{Dérivabilité en un point}
Soit $f:I\to\R$ et $a\in I$.

On dit que $f$ est dérivable en $a$ si la limite
\[
\lim_{x\to a,\ x\neq a}\frac{f(x)-f(a)}{x-a}
\]
existe et est finie.

Cette limite est appelée nombre dérivé de $f$ en $a$ et se note
\[
f'(a).
\]

De manière équivalente :
\[
f'(a)=\lim_{h\to0}\frac{f(a+h)-f(a)}{h}.
\]
\end{definitionbox}

\begin{definitionbox}{Dérivées latérales}
La dérivée à droite de $f$ en $a$ est
\[
f'_d(a)=\lim_{h\to0^+}\frac{f(a+h)-f(a)}{h},
\]
si cette limite existe.

La dérivée à gauche est
\[
f'_g(a)=\lim_{h\to0^-}\frac{f(a+h)-f(a)}{h},
\]
si cette limite existe.
\end{definitionbox}

\begin{theorembox}{Critère par dérivées latérales}
La fonction $f$ est dérivable en $a$ si et seulement si les dérivées à gauche et à droite existent et sont égales :
\[
f'_g(a)=f'_d(a).
\]
Dans ce cas,
\[
f'(a)=f'_g(a)=f'_d(a).
\]
\end{theorembox}

\begin{proofbox}{La fonction valeur absolue n'est pas dérivable en $0$}
Soit
\[
f(x)=|x|.
\]
On calcule
\[
\frac{f(h)-f(0)}{h}=\frac{|h|}{h}.
\]

Si $h>0$, alors $|h|=h$, donc
\[
\frac{|h|}{h}=1.
\]
Ainsi
\[
f'_d(0)=1.
\]

Si $h<0$, alors $|h|=-h$, donc
\[
\frac{|h|}{h}=-1.
\]
Ainsi
\[
f'_g(0)=-1.
\]

Les dérivées latérales sont différentes, donc $f$ n'est pas dérivable en $0$.
\end{proofbox}

%====================================================
\section{Dérivée et approximation locale}

\begin{theorembox}{Approximation affine}
Soit $f:I\to\R$ et $a\in I$.

La fonction $f$ est dérivable en $a$ si et seulement s'il existe une fonction $\varepsilon$ telle que
\[
\varepsilon(h)\xrightarrow[h\to0]{}0
\]
et
\[
f(a+h)=f(a)+hf'(a)+h\varepsilon(h).
\]

On écrit aussi :
\[
f(a+h)=f(a)+hf'(a)+o(h).
\]
\end{theorembox}

\begin{proofbox}{Dérivabilité équivaut à l'approximation affine}
Supposons d'abord que $f$ soit dérivable en $a$.

Pour $h\neq0$, posons
\[
\varepsilon(h)=\frac{f(a+h)-f(a)}{h}-f'(a).
\]
Par définition de la dérivée,
\[
\varepsilon(h)\to0.
\]
Donc
\[
\frac{f(a+h)-f(a)}{h}=f'(a)+\varepsilon(h),
\]
puis
\[
f(a+h)=f(a)+hf'(a)+h\varepsilon(h).
\]

Réciproquement, supposons qu'il existe $L\in\R$ et une fonction $\varepsilon$ telle que
\[
f(a+h)=f(a)+Lh+h\varepsilon(h),
\qquad \varepsilon(h)\to0.
\]
Alors, pour $h\neq0$,
\[
\frac{f(a+h)-f(a)}{h}=L+\varepsilon(h).
\]
En passant à la limite lorsque $h\to0$, on obtient
\[
\lim_{h\to0}\frac{f(a+h)-f(a)}{h}=L.
\]
Donc $f$ est dérivable en $a$ et
\[
f'(a)=L.
\]
\end{proofbox}

\begin{theorembox}{Dérivabilité implique continuité}
Si $f$ est dérivable en $a$, alors $f$ est continue en $a$.
\end{theorembox}

\begin{proofbox}{Dérivabilité implique continuité}
Si $f$ est dérivable en $a$, alors
\[
f(a+h)=f(a)+hf'(a)+h\varepsilon(h),
\qquad \varepsilon(h)\to0.
\]
Donc
\[
f(a+h)-f(a)=hf'(a)+h\varepsilon(h).
\]
Lorsque $h\to0$, on a
\[
hf'(a)\to0
\]
et
\[
h\varepsilon(h)\to0.
\]
Ainsi
\[
f(a+h)-f(a)\to0.
\]
Donc
\[
f(a+h)\to f(a).
\]
La fonction $f$ est continue en $a$.
\end{proofbox}

\begin{warningbox}{Réciproque fausse}
Une fonction peut être continue sans être dérivable.

Exemple :
\[
f(x)=|x|
\]
est continue en $0$, mais n'est pas dérivable en $0$.
\end{warningbox}

\begin{definitionbox}{Tangente}
Si $f$ est dérivable en $a$, la tangente à la courbe de $f$ au point d'abscisse $a$ a pour équation
\[
y=f(a)+f'(a)(x-a).
\]
\end{definitionbox}

%====================================================
\section{Règles de dérivation}

\begin{theorembox}{Linéarité}
Si $f$ et $g$ sont dérivables en $a$ et si $\lambda,\mu\in\R$, alors
\[
\lambda f+\mu g
\]
est dérivable en $a$ et
\[
(\lambda f+\mu g)'(a)=\lambda f'(a)+\mu g'(a).
\]
\end{theorembox}

\begin{proofbox}{Linéarité}
On écrit
\[
\frac{(\lambda f+\mu g)(a+h)-(\lambda f+\mu g)(a)}{h}
=
\lambda\frac{f(a+h)-f(a)}{h}
+
\mu\frac{g(a+h)-g(a)}{h}.
\]
En faisant tendre $h$ vers $0$, on obtient
\[
(\lambda f+\mu g)'(a)=\lambda f'(a)+\mu g'(a).
\]
\end{proofbox}

\begin{theorembox}{Produit}
Si $f$ et $g$ sont dérivables en $a$, alors $fg$ est dérivable en $a$ et
\[
(fg)'(a)=f'(a)g(a)+f(a)g'(a).
\]
\end{theorembox}

\begin{proofbox}{Dérivée d'un produit}
On part de
\[
f(a+h)g(a+h)-f(a)g(a).
\]
On ajoute et retranche $f(a)g(a+h)$ :
\[
f(a+h)g(a+h)-f(a)g(a+h)+f(a)g(a+h)-f(a)g(a).
\]
Donc
\[
f(a+h)g(a+h)-f(a)g(a)
=
\bigl(f(a+h)-f(a)\bigr)g(a+h)
+
f(a)\bigl(g(a+h)-g(a)\bigr).
\]
En divisant par $h$ :
\[
\frac{f(a+h)g(a+h)-f(a)g(a)}{h}
=
\frac{f(a+h)-f(a)}{h}g(a+h)
+
f(a)\frac{g(a+h)-g(a)}{h}.
\]
Comme $g$ est dérivable en $a$, elle est continue en $a$, donc
\[
g(a+h)\to g(a).
\]
En passant à la limite :
\[
(fg)'(a)=f'(a)g(a)+f(a)g'(a).
\]
\end{proofbox}

\begin{theorembox}{Inverse}
Si $g$ est dérivable en $a$ et si $g(a)\neq0$, alors $\dfrac1g$ est dérivable en $a$ et
\[
\left(\frac1g\right)'(a)=-\frac{g'(a)}{g(a)^2}.
\]
\end{theorembox}

\begin{proofbox}{Dérivée de l'inverse}
On écrit
\[
\frac{\frac1{g(a+h)}-\frac1{g(a)}}{h}
=
\frac{g(a)-g(a+h)}{h\,g(a+h)g(a)}.
\]
Donc
\[
\frac{\frac1{g(a+h)}-\frac1{g(a)}}{h}
=
-\frac{g(a+h)-g(a)}{h}\cdot \frac1{g(a+h)g(a)}.
\]
Comme $g$ est dérivable en $a$, elle est continue en $a$, donc
\[
g(a+h)\to g(a).
\]
Puisque $g(a)\neq0$,
\[
g(a+h)g(a)\to g(a)^2\neq0.
\]
Ainsi
\[
\left(\frac1g\right)'(a)
=
-\frac{g'(a)}{g(a)^2}.
\]
\end{proofbox}

\begin{theorembox}{Quotient}
Si $f$ et $g$ sont dérivables en $a$ et si $g(a)\neq0$, alors $\dfrac fg$ est dérivable en $a$ et
\[
\left(\frac fg\right)'(a)=\frac{f'(a)g(a)-f(a)g'(a)}{g(a)^2}.
\]
\end{theorembox}

\begin{proofbox}{Dérivée d'un quotient}
On écrit
\[
\frac fg=f\cdot \frac1g.
\]
D'après la formule du produit :
\[
\left(\frac fg\right)'(a)
=
f'(a)\frac1{g(a)}
+
f(a)\left(\frac1g\right)'(a).
\]
Avec la formule de l'inverse :
\[
\left(\frac1g\right)'(a)=-\frac{g'(a)}{g(a)^2}.
\]
Donc
\[
\left(\frac fg\right)'(a)
=
\frac{f'(a)}{g(a)}
-
\frac{f(a)g'(a)}{g(a)^2}.
\]
En mettant au même dénominateur :
\[
\left(\frac fg\right)'(a)
=
\frac{f'(a)g(a)-f(a)g'(a)}{g(a)^2}.
\]
\end{proofbox}

\begin{theorembox}{Composition}
Soient $g:I\to J$ et $f:J\to\R$.

Si $g$ est dérivable en $a$ et si $f$ est dérivable en $g(a)$, alors $f\circ g$ est dérivable en $a$ et
\[
(f\circ g)'(a)=f'(g(a))g'(a).
\]
\end{theorembox}

\begin{proofbox}{Dérivée d'une composée}
Comme $f$ est dérivable en $g(a)$, on peut écrire
\[
f(g(a)+k)=f(g(a))+kf'(g(a))+k\varepsilon(k),
\]
avec
\[
\varepsilon(k)\to0.
\]
Posons
\[
k=g(a+h)-g(a).
\]
Comme $g$ est dérivable en $a$, elle est continue en $a$, donc
\[
k\to0.
\]
Ainsi
\[
f(g(a+h))-f(g(a))
=
\bigl(g(a+h)-g(a)\bigr)f'(g(a))
+
\bigl(g(a+h)-g(a)\bigr)\varepsilon(g(a+h)-g(a)).
\]
En divisant par $h$ :
\[
\frac{f(g(a+h))-f(g(a))}{h}
=
\frac{g(a+h)-g(a)}{h}f'(g(a))
+
\frac{g(a+h)-g(a)}{h}\varepsilon(g(a+h)-g(a)).
\]
En passant à la limite :
\[
(f\circ g)'(a)=g'(a)f'(g(a)).
\]
\end{proofbox}

\begin{theorembox}{Réciproque}
Soit $f:I\to J$ une bijection continue strictement monotone.

Si $f$ est dérivable en $a\in I$ et si
\[
f'(a)\neq0,
\]
alors $f^{-1}$ est dérivable en $b=f(a)$ et
\[
(f^{-1})'(b)=\frac1{f'(a)}.
\]

Sous forme fonctionnelle :
\[
(f^{-1})'(y)=\frac1{f'(f^{-1}(y))}.
\]
\end{theorembox}

\begin{proofbox}{Dérivée d'une fonction réciproque}
Posons
\[
b=f(a).
\]
Pour $y\neq b$, on étudie
\[
\frac{f^{-1}(y)-f^{-1}(b)}{y-b}.
\]
Posons
\[
x=f^{-1}(y).
\]
Alors
\[
y=f(x)
\]
et
\[
f^{-1}(b)=a.
\]
Le quotient devient
\[
\frac{x-a}{f(x)-f(a)}
=
\frac1{\frac{f(x)-f(a)}{x-a}}.
\]
Lorsque $y\to b$, par continuité de $f^{-1}$, on a
\[
x=f^{-1}(y)\to f^{-1}(b)=a.
\]
Donc
\[
\frac{f(x)-f(a)}{x-a}\to f'(a).
\]
Comme $f'(a)\neq0$, on obtient
\[
\frac{x-a}{f(x)-f(a)}\to \frac1{f'(a)}.
\]
Ainsi
\[
(f^{-1})'(b)=\frac1{f'(a)}.
\]
\end{proofbox}

\begin{warningbox}{Erreur classique}
On ne doit pas écrire
\[
(f^{-1})'(x)=\frac1{f'(x)}.
\]
La bonne formule est
\[
(f^{-1})'(x)=\frac1{f'(f^{-1}(x))}.
\]
\end{warningbox}

%====================================================
\section{Théorèmes fondamentaux de la dérivation}

\begin{definitionbox}{Extremum local}
Soit $f:I\to\R$ et $a\in I$.

On dit que $f$ admet un maximum local en $a$ s'il existe $\eta>0$ tel que
\[
\forall x\in I,\quad |x-a|\leq\eta \Longrightarrow f(x)\leq f(a).
\]

On dit que $f$ admet un minimum local en $a$ s'il existe $\eta>0$ tel que
\[
\forall x\in I,\quad |x-a|\leq\eta \Longrightarrow f(a)\leq f(x).
\]
\end{definitionbox}

\begin{theorembox}{Condition nécessaire d'extremum local}
Soit $f:I\to\R$ dérivable en un point $a$ intérieur à $I$.

Si $f$ admet un extremum local en $a$, alors
\[
f'(a)=0.
\]
\end{theorembox}

\begin{proofbox}{Condition nécessaire d'extremum}
Supposons que $f$ admette un maximum local en $a$.

Il existe $\eta>0$ tel que, pour $|h|$ assez petit,
\[
f(a+h)\leq f(a).
\]

Si $h>0$, alors
\[
\frac{f(a+h)-f(a)}{h}\leq0.
\]
En passant à la limite quand $h\to0^+$ :
\[
f'_d(a)\leq0.
\]

Si $h<0$, alors le numérateur est encore négatif ou nul, mais le dénominateur est négatif, donc
\[
\frac{f(a+h)-f(a)}{h}\geq0.
\]
En passant à la limite quand $h\to0^-$ :
\[
f'_g(a)\geq0.
\]

Comme $f$ est dérivable en $a$,
\[
f'_g(a)=f'_d(a)=f'(a).
\]
Donc
\[
f'(a)\leq0
\quad\text{et}\quad
f'(a)\geq0.
\]
Ainsi
\[
f'(a)=0.
\]

Le cas d'un minimum local est analogue.
\end{proofbox}

\begin{warningbox}{Point critique}
Un point critique n'est pas forcément un extremum.

Exemple :
\[
f(x)=x^3.
\]
On a
\[
f'(0)=0,
\]
mais $0$ n'est ni un maximum local ni un minimum local.
\end{warningbox}

\begin{theorembox}{Théorème de Rolle}
Soit $f:[a,b]\to\R$.

Si :
\begin{enumerate}[label=\arabic*)]
    \item $f$ est continue sur $[a,b]$ ;
    \item $f$ est dérivable sur $]a,b[$ ;
    \item $f(a)=f(b)$ ;
\end{enumerate}
alors il existe $c\in]a,b[$ tel que
\[
f'(c)=0.
\]
\end{theorembox}

\begin{proofbox}{Théorème de Rolle}
Comme $f$ est continue sur le segment $[a,b]$, elle atteint son minimum et son maximum.

Il existe donc $\alpha,\beta\in[a,b]$ tels que
\[
f(\alpha)=\min_{[a,b]} f,
\qquad
f(\beta)=\max_{[a,b]} f.
\]

Si $f$ est constante sur $[a,b]$, alors
\[
f'(x)=0
\]
pour tout $x\in]a,b[$, et le résultat est immédiat.

Sinon, le minimum et le maximum sont distincts. Comme
\[
f(a)=f(b),
\]
au moins l'un des extrema est atteint en un point intérieur $c\in]a,b[$.

En ce point $c$, la fonction admet un extremum local et est dérivable. Donc, par la condition nécessaire d'extremum :
\[
f'(c)=0.
\]
\end{proofbox}

\begin{theorembox}{Théorème des accroissements finis}
Soit $f:[a,b]\to\R$.

Si :
\begin{enumerate}[label=\arabic*)]
    \item $f$ est continue sur $[a,b]$ ;
    \item $f$ est dérivable sur $]a,b[$ ;
\end{enumerate}
alors il existe $c\in]a,b[$ tel que
\[
f(b)-f(a)=f'(c)(b-a).
\]
Autrement dit :
\[
\frac{f(b)-f(a)}{b-a}=f'(c).
\]
\end{theorembox}

\begin{proofbox}{Théorème des accroissements finis}
On considère la fonction auxiliaire
\[
g(x)=f(x)-\frac{f(b)-f(a)}{b-a}(x-a).
\]
La fonction $g$ est continue sur $[a,b]$ et dérivable sur $]a,b[$.

On calcule :
\[
g(a)=f(a),
\]
et
\[
g(b)=f(b)-\frac{f(b)-f(a)}{b-a}(b-a)=f(a).
\]
Donc
\[
g(a)=g(b).
\]

Par le théorème de Rolle, il existe $c\in]a,b[$ tel que
\[
g'(c)=0.
\]
Or
\[
g'(x)=f'(x)-\frac{f(b)-f(a)}{b-a}.
\]
Donc
\[
f'(c)=\frac{f(b)-f(a)}{b-a}.
\]
\end{proofbox}

\begin{theorembox}{Variations et signe de la dérivée}
Soit $f:I\to\R$ dérivable sur un intervalle $I$.

\begin{enumerate}[label=\arabic*)]
    \item Si $f'(x)\geq0$ pour tout $x\in I$, alors $f$ est croissante sur $I$.
    \item Si $f'(x)\leq0$ pour tout $x\in I$, alors $f$ est décroissante sur $I$.
    \item Si $f'(x)>0$ pour tout $x\in I$, alors $f$ est strictement croissante sur $I$.
    \item Si $f'(x)<0$ pour tout $x\in I$, alors $f$ est strictement décroissante sur $I$.
\end{enumerate}
\end{theorembox}

\begin{proofbox}{Signe de la dérivée et variations}
Soient $x,y\in I$ avec $x<y$.

Par le théorème des accroissements finis appliqué à $f$ sur $[x,y]$, il existe $c\in]x,y[$ tel que
\[
f(y)-f(x)=f'(c)(y-x).
\]

Comme
\[
y-x>0,
\]
le signe de $f(y)-f(x)$ est celui de $f'(c)$.

On en déduit les résultats de croissance ou décroissance.
\end{proofbox}

\begin{theorembox}{Inégalité des accroissements finis}
Soit $f:I\to\R$ dérivable sur un intervalle $I$.

S'il existe $K\geq0$ tel que
\[
\forall x\in I,\quad |f'(x)|\leq K,
\]
alors
\[
\forall x,y\in I,\quad |f(x)-f(y)|\leq K|x-y|.
\]
\end{theorembox}

\begin{proofbox}{Inégalité des accroissements finis}
Soient $x,y\in I$.

Si $x=y$, le résultat est évident.

Supposons $x<y$. Par le théorème des accroissements finis, il existe $c\in]x,y[$ tel que
\[
f(y)-f(x)=f'(c)(y-x).
\]
Donc
\[
|f(y)-f(x)|=|f'(c)||y-x|.
\]
Comme
\[
|f'(c)|\leq K,
\]
on obtient
\[
|f(y)-f(x)|\leq K|y-x|.
\]
Le cas $y<x$ est identique en échangeant $x$ et $y$.
\end{proofbox}

%====================================================
\section{Table complète des fonctions usuelles}

\subsection{Fonctions puissances entières}

\begin{theorembox}{Puissances entières}
Pour tout $n\in\N$, la fonction
\[
x\mapsto x^n
\]
est dérivable sur $\R$ et, pour $n\geq1$,
\[
(x^n)'=nx^{n-1}.
\]
\end{theorembox}

\begin{proofbox}{Dérivée de $x^n$}
Soit
\[
f(x)=x^n.
\]
Pour $h\neq0$ :
\[
\frac{(x+h)^n-x^n}{h}.
\]
Par la formule du binôme :
\[
(x+h)^n=\sum_{k=0}^n \binom{n}{k}x^{n-k}h^k.
\]
Donc
\[
(x+h)^n-x^n
=
nx^{n-1}h+\sum_{k=2}^n \binom{n}{k}x^{n-k}h^k.
\]
En divisant par $h$ :
\[
\frac{(x+h)^n-x^n}{h}
=
nx^{n-1}
+
\sum_{k=2}^n \binom{n}{k}x^{n-k}h^{k-1}.
\]
Lorsque $h\to0$, tous les termes de la somme tendent vers $0$.

Ainsi
\[
(x^n)'=nx^{n-1}.
\]
\end{proofbox}

\subsection{Fonction inverse}

\begin{theorembox}{Inverse}
La fonction
\[
x\mapsto \frac1x
\]
est dérivable sur $\R^\ast$ et
\[
\left(\frac1x\right)'=-\frac1{x^2}.
\]
\end{theorembox}

\begin{proofbox}{Dérivée de l'inverse}
Soit $f(x)=1/x$.

Pour $h\neq0$ et $x+h\neq0$ :
\[
\frac{f(x+h)-f(x)}{h}
=
\frac{\frac1{x+h}-\frac1x}{h}.
\]
On simplifie :
\[
\frac{\frac{x-(x+h)}{x(x+h)}}{h}
=
\frac{-h}{h x(x+h)}
=
-\frac1{x(x+h)}.
\]
Lorsque $h\to0$ :
\[
-\frac1{x(x+h)}\to -\frac1{x^2}.
\]
Donc
\[
\left(\frac1x\right)'=-\frac1{x^2}.
\]
\end{proofbox}

\subsection{Racine carrée}

\begin{theorembox}{Racine carrée}
La fonction
\[
x\mapsto \sqrt{x}
\]
est dérivable sur $]0,+\infty[$ et
\[
(\sqrt{x})'=\frac1{2\sqrt{x}}.
\]

Elle est continue sur $[0,+\infty[$, mais elle n'est pas dérivable en $0$.
\end{theorembox}

\begin{proofbox}{Dérivée de la racine carrée}
Pour $x>0$ et $h$ assez petit :
\[
\frac{\sqrt{x+h}-\sqrt{x}}{h}.
\]
On multiplie par la quantité conjuguée :
\[
\frac{\sqrt{x+h}-\sqrt{x}}{h}
\cdot
\frac{\sqrt{x+h}+\sqrt{x}}{\sqrt{x+h}+\sqrt{x}}
=
\frac{x+h-x}{h(\sqrt{x+h}+\sqrt{x})}.
\]
Donc
\[
\frac{\sqrt{x+h}-\sqrt{x}}{h}
=
\frac1{\sqrt{x+h}+\sqrt{x}}.
\]
En faisant tendre $h$ vers $0$ :
\[
(\sqrt{x})'=\frac1{2\sqrt{x}}.
\]
\end{proofbox}

\subsection{Logarithme népérien}

\begin{definitionbox}{Logarithme}
La fonction logarithme népérien est définie sur $]0,+\infty[$.

Elle vérifie :
\[
\ln(1)=0,
\]
\[
\ln(xy)=\ln x+\ln y,
\]
et
\[
(\ln x)'=\frac1x.
\]
\end{definitionbox}

\begin{theorembox}{Propriétés du logarithme}
Pour tous $x,y>0$ :
\[
\ln(xy)=\ln x+\ln y,
\]
\[
\ln\left(\frac xy\right)=\ln x-\ln y,
\]
\[
\ln(x^r)=r\ln x
\]
pour tout réel $r$ pour lequel l'expression a un sens.

De plus :
\[
\lim_{x\to0^+}\ln x=-\infty,
\qquad
\lim_{x\to+\infty}\ln x=+\infty.
\]
\end{theorembox}

\begin{theorembox}{Dérivée de $\ln u$}
Si $u$ est dérivable et strictement positive sur un intervalle $I$, alors
\[
(\ln u)'=\frac{u'}u.
\]
\end{theorembox}

\begin{proofbox}{Dérivée de $\ln u$}
La fonction $\ln u$ est une composée :
\[
x\mapsto u(x)\mapsto \ln(u(x)).
\]
Comme
\[
(\ln t)'=\frac1t,
\]
on applique la formule de dérivation d'une composée :
\[
(\ln u)'(x)=\ln'(u(x))u'(x)=\frac{u'(x)}{u(x)}.
\]
\end{proofbox}

\subsection{Exponentielle}

\begin{definitionbox}{Exponentielle}
La fonction exponentielle est la réciproque de la fonction logarithme :
\[
\exp : \R\to]0,+\infty[.
\]
On note
\[
\exp(x)=\e^x.
\]

On a :
\[
\ln(\e^x)=x,
\qquad
\e^{\ln x}=x\quad(x>0).
\]
\end{definitionbox}

\begin{theorembox}{Propriétés de l'exponentielle}
Pour tous $x,y\in\R$ :
\[
\e^{x+y}=\e^x\e^y,
\]
\[
\e^{-x}=\frac1{\e^x},
\]
\[
\e^0=1.
\]

De plus :
\[
(\e^x)'=\e^x,
\]
\[
\lim_{x\to-\infty}\e^x=0,
\qquad
\lim_{x\to+\infty}\e^x=+\infty.
\]
\end{theorembox}

\begin{proofbox}{Dérivée de l'exponentielle}
Comme $\exp$ est la réciproque de $\ln$, on utilise la formule de dérivation d'une réciproque.

Soit
\[
y=\e^x.
\]
Alors
\[
x=\ln y.
\]
On sait que
\[
(\ln)'(y)=\frac1y.
\]
Donc
\[
(\exp)'(x)
=
\frac1{\ln'(\exp x)}
=
\frac1{1/\e^x}
=
\e^x.
\]
\end{proofbox}

\begin{theorembox}{Dérivée de $\e^u$}
Si $u$ est dérivable sur $I$, alors
\[
(\e^u)'=u'\e^u.
\]
\end{theorembox}

\subsection{Puissances réelles}

\begin{definitionbox}{Puissance réelle}
Pour $a\in\R$ et $x>0$, on définit
\[
x^a=\e^{a\ln x}.
\]
\end{definitionbox}

\begin{theorembox}{Dérivée de $x^a$}
Pour tout réel $a$, la fonction
\[
x\mapsto x^a
\]
est dérivable sur $]0,+\infty[$ et
\[
(x^a)'=ax^{a-1}.
\]
\end{theorembox}

\begin{proofbox}{Dérivée de $x^a$}
On écrit
\[
x^a=\e^{a\ln x}.
\]
Alors, par dérivation d'une composée :
\[
(x^a)'=
\left(a\frac1x\right)\e^{a\ln x}.
\]
Donc
\[
(x^a)'=\frac ax x^a=ax^{a-1}.
\]
\end{proofbox}

%====================================================
\section{Fonctions trigonométriques}

\begin{theorembox}{Sinus et cosinus}
Les fonctions $\sin$ et $\cos$ sont dérivables sur $\R$ et
\[
(\sin x)'=\cos x,
\]
\[
(\cos x)'=-\sin x.
\]

Elles vérifient :
\[
\cos^2 x+\sin^2 x=1.
\]
\end{theorembox}

\begin{theorembox}{Tangente}
La fonction tangente est définie par
\[
\tan x=\frac{\sin x}{\cos x}
\]
sur
\[
\R\setminus\left\{\frac\pi2+k\pi,\ k\in\mathbb{Z}\right\}.
\]

Elle est dérivable sur chaque intervalle de son domaine et
\[
(\tan x)'=1+\tan^2 x=\frac1{\cos^2 x}.
\]
\end{theorembox}

\begin{proofbox}{Dérivée de la tangente}
On utilise la formule du quotient :
\[
(\tan x)'=
\left(\frac{\sin x}{\cos x}\right)'.
\]
Donc
\[
(\tan x)'
=
\frac{\cos x\cos x-\sin x(-\sin x)}{\cos^2 x}.
\]
Ainsi
\[
(\tan x)'=
\frac{\cos^2 x+\sin^2 x}{\cos^2 x}
=
\frac1{\cos^2 x}.
\]
Comme
\[
1+\tan^2 x
=
1+\frac{\sin^2 x}{\cos^2 x}
=
\frac{\cos^2 x+\sin^2 x}{\cos^2 x}
=
\frac1{\cos^2 x},
\]
on obtient aussi
\[
(\tan x)'=1+\tan^2 x.
\]
\end{proofbox}

\begin{methodbox}{Dériver une composée trigonométrique}
Si $u$ est dérivable :
\[
(\sin u)'=u'\cos u,
\]
\[
(\cos u)'=-u'\sin u,
\]
\[
(\tan u)'=\frac{u'}{\cos^2 u}=u'(1+\tan^2 u).
\]
\end{methodbox}

%====================================================
\section{Fonctions trigonométriques réciproques}

\subsection{Arcsinus}

\begin{definitionbox}{Arcsinus}
La fonction sinus réalise une bijection de
\[
\left[-\frac\pi2,\frac\pi2\right]
\]
sur
\[
[-1,1].
\]

Sa réciproque est appelée arcsinus :
\[
\arcsin : [-1,1]\to\left[-\frac\pi2,\frac\pi2\right].
\]
\end{definitionbox}

\begin{theorembox}{Dérivée de l'arcsinus}
La fonction $\arcsin$ est dérivable sur $]-1,1[$ et
\[
(\arcsin x)'=\frac1{\sqrt{1-x^2}}.
\]
\end{theorembox}

\begin{proofbox}{Dérivée de l'arcsinus}
Soit
\[
y=\arcsin x.
\]
Alors
\[
x=\sin y,
\]
avec
\[
y\in\left]-\frac\pi2,\frac\pi2\right[.
\]
Par dérivation d'une réciproque :
\[
(\arcsin)'(x)=\frac1{\cos y}.
\]
Or sur
\[
\left]-\frac\pi2,\frac\pi2\right[,
\]
on a
\[
\cos y>0.
\]
De plus
\[
\cos^2 y=1-\sin^2 y=1-x^2.
\]
Donc
\[
\cos y=\sqrt{1-x^2}.
\]
Ainsi
\[
(\arcsin x)'=\frac1{\sqrt{1-x^2}}.
\]
\end{proofbox}

\subsection{Arccosinus}

\begin{definitionbox}{Arccosinus}
La fonction cosinus réalise une bijection de
\[
[0,\pi]
\]
sur
\[
[-1,1].
\]

Sa réciproque est appelée arccosinus :
\[
\arccos : [-1,1]\to[0,\pi].
\]
\end{definitionbox}

\begin{theorembox}{Dérivée de l'arccosinus}
La fonction $\arccos$ est dérivable sur $]-1,1[$ et
\[
(\arccos x)'=-\frac1{\sqrt{1-x^2}}.
\]
\end{theorembox}

\begin{proofbox}{Dérivée de l'arccosinus}
Soit
\[
y=\arccos x.
\]
Alors
\[
x=\cos y,
\]
avec
\[
y\in]0,\pi[.
\]
Par dérivation d'une réciproque :
\[
(\arccos)'(x)=\frac1{-\sin y}.
\]
Or sur $]0,\pi[$ :
\[
\sin y>0.
\]
Comme
\[
\sin^2 y=1-\cos^2 y=1-x^2,
\]
on a
\[
\sin y=\sqrt{1-x^2}.
\]
Donc
\[
(\arccos x)'=-\frac1{\sqrt{1-x^2}}.
\]
\end{proofbox}

\subsection{Arctangente}

\begin{definitionbox}{Arctangente}
La fonction tangente réalise une bijection de
\[
\left]-\frac\pi2,\frac\pi2\right[
\]
sur $\R$.

Sa réciproque est appelée arctangente :
\[
\arctan:\R\to\left]-\frac\pi2,\frac\pi2\right[.
\]
\end{definitionbox}

\begin{theorembox}{Dérivée de l'arctangente}
La fonction $\arctan$ est dérivable sur $\R$ et
\[
(\arctan x)'=\frac1{1+x^2}.
\]
\end{theorembox}

\begin{proofbox}{Dérivée de l'arctangente}
Soit
\[
y=\arctan x.
\]
Alors
\[
x=\tan y.
\]
Par dérivation d'une réciproque :
\[
(\arctan)'(x)=\frac1{1+\tan^2 y}.
\]
Comme
\[
\tan y=x,
\]
on obtient
\[
(\arctan x)'=\frac1{1+x^2}.
\]
\end{proofbox}

\begin{theorembox}{Relation entre arcsinus et arccosinus}
Pour tout $x\in[-1,1]$ :
\[
\arcsin x+\arccos x=\frac\pi2.
\]
\end{theorembox}

\begin{proofbox}{Relation $\arcsin+\arccos$}
Posons
\[
F(x)=\arcsin x+\arccos x.
\]
Sur $]-1,1[$ :
\[
F'(x)=\frac1{\sqrt{1-x^2}}-\frac1{\sqrt{1-x^2}}=0.
\]
Donc $F$ est constante sur $]-1,1[$.

En évaluant en $0$ :
\[
F(0)=\arcsin 0+\arccos 0=0+\frac\pi2=\frac\pi2.
\]
Par continuité des deux fonctions aux bornes, l'égalité reste vraie sur $[-1,1]$.
\end{proofbox}

%====================================================
\section{Fonctions hyperboliques}

\begin{definitionbox}{Fonctions hyperboliques}
Pour tout $x\in\R$, on définit
\[
\ch x=\frac{\e^x+\e^{-x}}2,
\]
\[
\sh x=\frac{\e^x-\e^{-x}}2,
\]
\[
\thh x=\frac{\sh x}{\ch x}.
\]
\end{definitionbox}

\begin{theorembox}{Identité fondamentale}
Pour tout $x\in\R$ :
\[
\ch^2 x-\sh^2 x=1.
\]
\end{theorembox}

\begin{proofbox}{Identité hyperbolique}
On calcule :
\[
\ch^2 x-\sh^2 x
=
\left(\frac{\e^x+\e^{-x}}2\right)^2
-
\left(\frac{\e^x-\e^{-x}}2\right)^2.
\]
On utilise
\[
(a+b)^2-(a-b)^2=4ab.
\]
Avec
\[
a=\e^x,\qquad b=\e^{-x},
\]
on obtient
\[
\ch^2 x-\sh^2 x
=
\frac{4\e^x\e^{-x}}4
=
1.
\]
\end{proofbox}

\begin{theorembox}{Dérivées hyperboliques}
Les fonctions $\ch$, $\sh$ et $\thh$ sont dérivables sur $\R$ et
\[
(\ch x)'=\sh x,
\]
\[
(\sh x)'=\ch x,
\]
\[
(\thh x)'=1-\thh^2 x=\frac1{\ch^2 x}.
\]
\end{theorembox}

\begin{proofbox}{Dérivées hyperboliques}
On dérive les définitions :
\[
(\ch x)'
=
\left(\frac{\e^x+\e^{-x}}2\right)'
=
\frac{\e^x-\e^{-x}}2
=
\sh x.
\]
De même :
\[
(\sh x)'
=
\left(\frac{\e^x-\e^{-x}}2\right)'
=
\frac{\e^x+\e^{-x}}2
=
\ch x.
\]

Pour $\thh$ :
\[
\thh x=\frac{\sh x}{\ch x}.
\]
Donc
\[
(\thh x)'
=
\frac{\ch x\cdot \ch x-\sh x\cdot \sh x}{\ch^2 x}.
\]
Ainsi
\[
(\thh x)'=
\frac{\ch^2 x-\sh^2 x}{\ch^2 x}
=
\frac1{\ch^2 x}.
\]
Enfin
\[
1-\thh^2 x
=
1-\frac{\sh^2 x}{\ch^2 x}
=
\frac{\ch^2 x-\sh^2 x}{\ch^2 x}
=
\frac1{\ch^2 x}.
\]
\end{proofbox}

%====================================================
\section{Tableau récapitulatif des dérivées}

\renewcommand{\arraystretch}{1.6}
\begin{center}
\begin{tabular}{|>{\centering\arraybackslash}m{5cm}|>{\centering\arraybackslash}m{5cm}|>{\centering\arraybackslash}m{4cm}|}
\hline
\textbf{Fonction} & \textbf{Dérivée} & \textbf{Domaine} \\
\hline
$c$ & $0$ & $\R$ \\
\hline
$x$ & $1$ & $\R$ \\
\hline
$x^n$ & $nx^{n-1}$ & $\R$ \\
\hline
$\dfrac1x$ & $-\dfrac1{x^2}$ & $\R^\ast$ \\
\hline
$\sqrt{x}$ & $\dfrac1{2\sqrt{x}}$ & $]0,+\infty[$ \\
\hline
$\ln x$ & $\dfrac1x$ & $]0,+\infty[$ \\
\hline
$\e^x$ & $\e^x$ & $\R$ \\
\hline
$x^\alpha$ & $\alpha x^{\alpha-1}$ & $]0,+\infty[$ \\
\hline
$\sin x$ & $\cos x$ & $\R$ \\
\hline
$\cos x$ & $-\sin x$ & $\R$ \\
\hline
$\tan x$ & $1+\tan^2x=\dfrac1{\cos^2x}$ & $\R\setminus\{\frac\pi2+k\pi\}$ \\
\hline
$\arcsin x$ & $\dfrac1{\sqrt{1-x^2}}$ & $]-1,1[$ \\
\hline
$\arccos x$ & $-\dfrac1{\sqrt{1-x^2}}$ & $]-1,1[$ \\
\hline
$\arctan x$ & $\dfrac1{1+x^2}$ & $\R$ \\
\hline
$\ch x$ & $\sh x$ & $\R$ \\
\hline
$\sh x$ & $\ch x$ & $\R$ \\
\hline
$\thh x$ & $1-\thh^2x=\dfrac1{\ch^2x}$ & $\R$ \\
\hline
\end{tabular}
\end{center}

%====================================================
\section{Formules de dérivation composées}

\begin{center}
\renewcommand{\arraystretch}{1.7}
\begin{tabular}{|>{\centering\arraybackslash}m{6cm}|>{\centering\arraybackslash}m{7cm}|}
\hline
\textbf{Fonction} & \textbf{Dérivée} \\
\hline
$u+v$ & $u'+v'$ \\
\hline
$\lambda u$ & $\lambda u'$ \\
\hline
$uv$ & $u'v+uv'$ \\
\hline
$\dfrac1u$ & $-\dfrac{u'}{u^2}$ \\
\hline
$\dfrac uv$ & $\dfrac{u'v-uv'}{v^2}$ \\
\hline
$u^n$ & $nu'u^{n-1}$ \\
\hline
$\sqrt{u}$ & $\dfrac{u'}{2\sqrt{u}}$ \\
\hline
$\ln u$ & $\dfrac{u'}u$ \\
\hline
$\e^u$ & $u'\e^u$ \\
\hline
$u^\alpha$ & $\alpha u'u^{\alpha-1}$ \\
\hline
$\sin u$ & $u'\cos u$ \\
\hline
$\cos u$ & $-u'\sin u$ \\
\hline
$\tan u$ & $\dfrac{u'}{\cos^2u}$ \\
\hline
$\arcsin u$ & $\dfrac{u'}{\sqrt{1-u^2}}$ \\
\hline
$\arccos u$ & $-\dfrac{u'}{\sqrt{1-u^2}}$ \\
\hline
$\arctan u$ & $\dfrac{u'}{1+u^2}$ \\
\hline
\end{tabular}
\end{center}

%====================================================
\section{Inégalités classiques à connaître}

\begin{theorembox}{Inégalités fondamentales}
Pour tout $x>-1$ :
\[
\ln(1+x)\leq x.
\]

Pour tout $x>0$ :
\[
\ln x\leq x-1.
\]

Pour tout $x\in\R$ :
\[
\e^x\geq 1+x.
\]

Pour tout $x\in\R$ :
\[
|\sin x|\leq |x|.
\]

Pour tout $x\in\R$ :
\[
|\sin x-\sin y|\leq |x-y|.
\]
\end{theorembox}

\begin{proofbox}{$\ln(1+x)\leq x$}
Soit
\[
F(x)=x-\ln(1+x),
\]
définie sur $]-1,+\infty[$.

On dérive :
\[
F'(x)=1-\frac1{1+x}=\frac{x}{1+x}.
\]
Comme $1+x>0$, le signe de $F'(x)$ est celui de $x$.

Donc $F$ est décroissante sur $]-1,0]$ et croissante sur $[0,+\infty[$.

Elle admet donc un minimum en $0$.

Or
\[
F(0)=0-\ln1=0.
\]
Ainsi
\[
F(x)\geq0.
\]
Donc
\[
\ln(1+x)\leq x.
\]
\end{proofbox}

\begin{proofbox}{$\e^x\geq 1+x$}
Soit
\[
F(x)=\e^x-1-x.
\]
On dérive :
\[
F'(x)=\e^x-1.
\]
Comme $\e^x-1<0$ si $x<0$, et $\e^x-1>0$ si $x>0$, la fonction $F$ est décroissante sur $]-\infty,0]$ puis croissante sur $[0,+\infty[$.

Elle admet donc un minimum en $0$.

Or
\[
F(0)=1-1-0=0.
\]
Donc
\[
F(x)\geq0.
\]
Ainsi
\[
\e^x\geq1+x.
\]
\end{proofbox}

\begin{proofbox}{$|\sin x-\sin y|\leq |x-y|$}
La fonction sinus est dérivable sur $\R$ et
\[
(\sin)'=\cos.
\]
Or
\[
|\cos t|\leq1
\]
pour tout $t\in\R$.

Par l'inégalité des accroissements finis, pour tous $x,y\in\R$ :
\[
|\sin x-\sin y|\leq |x-y|.
\]
\end{proofbox}

%====================================================
\section{Méthodes classiques en analyse}

\begin{methodbox}{Étudier une fonction}
Pour étudier une fonction $f$ :
\begin{enumerate}[label=\arabic*)]
    \item déterminer son domaine de définition ;
    \item étudier les symétries éventuelles : parité, périodicité ;
    \item étudier les limites aux bornes du domaine ;
    \item calculer la dérivée ;
    \item étudier le signe de la dérivée ;
    \item dresser le tableau de variations ;
    \item en déduire les extrema, images d'intervalles, solutions d'équations ;
    \item si besoin, étudier les asymptotes ou tangentes.
\end{enumerate}
\end{methodbox}

\begin{methodbox}{Montrer une inégalité}
Pour montrer une inégalité du type
\[
A(x)\leq B(x),
\]
on pose souvent
\[
F(x)=B(x)-A(x).
\]
Puis :
\begin{enumerate}[label=\arabic*)]
    \item on étudie $F'$ ;
    \item on trouve les variations de $F$ ;
    \item on montre que le minimum de $F$ est positif.
\end{enumerate}
\end{methodbox}

\begin{methodbox}{Montrer l'unicité d'une solution}
Pour montrer qu'une équation
\[
f(x)=0
\]
admet une unique solution sur un intervalle $I$ :
\begin{enumerate}[label=\arabic*)]
    \item existence : utiliser le théorème des valeurs intermédiaires ;
    \item unicité : montrer que $f$ est strictement monotone.
\end{enumerate}
\end{methodbox}

\begin{methodbox}{Utiliser Rolle}
On utilise Rolle pour obtenir une annulation de dérivée.

Il faut vérifier :
\begin{enumerate}[label=\arabic*)]
    \item continuité sur $[a,b]$ ;
    \item dérivabilité sur $]a,b[$ ;
    \item égalité $f(a)=f(b)$.
\end{enumerate}
Alors :
\[
\exists c\in]a,b[,\quad f'(c)=0.
\]
\end{methodbox}

\begin{methodbox}{Utiliser le TAF}
On utilise le théorème des accroissements finis pour relier :
\[
f(b)-f(a)
\]
à une valeur de la dérivée.

Si $f$ est continue sur $[a,b]$ et dérivable sur $]a,b[$, alors
\[
\exists c\in]a,b[,\quad f(b)-f(a)=f'(c)(b-a).
\]
\end{methodbox}

\begin{methodbox}{Utiliser l'IAF}
Pour majorer une différence :
\[
|f(x)-f(y)|,
\]
on cherche un majorant $K$ de $|f'|$.

Si
\[
|f'(t)|\leq K
\]
sur l'intervalle entre $x$ et $y$, alors
\[
|f(x)-f(y)|\leq K|x-y|.
\]
\end{methodbox}

%====================================================
\section{Démonstrations importantes à savoir refaire}

\subsection{Dérivée de la réciproque}

À savoir refaire parfaitement :
\[
(f^{-1})'(y)=\frac1{f'(f^{-1}(y))}.
\]

Idée :
\[
\frac{f^{-1}(y)-f^{-1}(b)}{y-b}
=
\frac{x-a}{f(x)-f(a)}
=
\frac1{\frac{f(x)-f(a)}{x-a}}.
\]

\subsection{Théorème de Rolle}

À savoir refaire parfaitement :
\begin{enumerate}[label=\arabic*)]
    \item continuité sur un segment ;
    \item existence d'un maximum et d'un minimum ;
    \item si la fonction n'est pas constante, l'un des extrema est atteint à l'intérieur ;
    \item condition nécessaire d'extremum : $f'(c)=0$.
\end{enumerate}

\subsection{Théorème des accroissements finis}

À savoir refaire parfaitement :
\[
g(x)=f(x)-\frac{f(b)-f(a)}{b-a}(x-a).
\]
Puis :
\[
g(a)=g(b),
\]
donc Rolle donne
\[
g'(c)=0.
\]
Ainsi :
\[
f'(c)=\frac{f(b)-f(a)}{b-a}.
\]

\subsection{Inégalité des accroissements finis}

À savoir refaire parfaitement :
\[
f(y)-f(x)=f'(c)(y-x)
\]
puis
\[
|f(y)-f(x)|=|f'(c)||y-x|\leq K|y-x|.
\]

\subsection{Dérivabilité implique continuité}

À savoir refaire parfaitement :
\[
f(a+h)=f(a)+hf'(a)+o(h),
\]
donc
\[
f(a+h)-f(a)=hf'(a)+o(h)\to0.
\]

%====================================================
\section{Erreurs classiques}

\begin{warningbox}{Erreur 1 : confondre continuité et dérivabilité}
La dérivabilité implique la continuité.

La continuité n'implique pas la dérivabilité.

Exemple :
\[
x\mapsto |x|
\]
est continue en $0$, mais pas dérivable en $0$.
\end{warningbox}

\begin{warningbox}{Erreur 2 : oublier le domaine}
Avant de dériver :
\begin{itemize}
    \item vérifier que la fonction est définie ;
    \item vérifier que le dénominateur ne s'annule pas ;
    \item vérifier que l'argument d'un logarithme est strictement positif ;
    \item vérifier que l'argument de $\arcsin$ ou $\arccos$ appartient à $[-1,1]$.
\end{itemize}
\end{warningbox}

\begin{warningbox}{Erreur 3 : mal dériver une composée}
On ne dérive pas
\[
\ln(u)
\]
comme
\[
\frac1u.
\]
La bonne formule est
\[
(\ln u)'=\frac{u'}u.
\]

De même :
\[
(\e^u)'=u'\e^u.
\]
\end{warningbox}

\begin{warningbox}{Erreur 4 : mauvaise formule pour la réciproque}
La formule correcte est
\[
(f^{-1})'(y)=\frac1{f'(f^{-1}(y))}.
\]
\end{warningbox}

\begin{warningbox}{Erreur 5 : point critique ne veut pas dire extremum}
Si $f'(a)=0$, alors $a$ est un point critique.

Mais il n'y a pas forcément d'extremum.

Exemple :
\[
f(x)=x^3,
\qquad
f'(0)=0,
\]
mais $0$ n'est pas un extremum.
\end{warningbox}

%====================================================
\section{Mini-fiche finale à mémoriser}

\begin{theorembox}{À connaître par coeur}
\[
f'(a)=\lim_{h\to0}\frac{f(a+h)-f(a)}h.
\]

\[
f(a+h)=f(a)+hf'(a)+o(h).
\]

\[
f\text{ dérivable en }a \Longrightarrow f\text{ continue en }a.
\]

\[
(fg)'=f'g+fg'.
\]

\[
\left(\frac fg\right)'=\frac{f'g-fg'}{g^2}.
\]

\[
(f\circ g)'=(f'\circ g)g'.
\]

\[
(f^{-1})'(y)=\frac1{f'(f^{-1}(y))}.
\]

\[
(\ln u)'=\frac{u'}u.
\]

\[
(\e^u)'=u'\e^u.
\]

\[
(\sin u)'=u'\cos u.
\]

\[
(\cos u)'=-u'\sin u.
\]

\[
(\arctan u)'=\frac{u'}{1+u^2}.
\]

\[
(\arcsin u)'=\frac{u'}{\sqrt{1-u^2}}.
\]

\[
(\arccos u)'=-\frac{u'}{\sqrt{1-u^2}}.
\]

Rolle :
\[
f(a)=f(b)\Longrightarrow \exists c,\ f'(c)=0.
\]

TAF :
\[
f(b)-f(a)=f'(c)(b-a).
\]

IAF :
\[
|f'|\leq K\Longrightarrow |f(x)-f(y)|\leq K|x-y|.
\]
\end{theorembox}

\end{document}